% from ega2-1-fr.tex
\oldpage[II]{147}
en outre $Z$ est projectif sur $Y$, et on voit qu'on peut se borner au cas où
$P$ est \emph{intègre} et $j$ \emph{dominant et birationnel}, et tout revient à
voir que $j$ est \emph{surjectif}. Or, soit $z\in P$~; $\mathcal{O}_z$ est un
anneau local intègre (resp. intègre et noethérien) dont le corps des fractions
est
\[
L=R(P)=R(X')=R(X).
\]
On peut se borner au cas où $z$ n'est pas le point générique de $P$. Il y a par
suite (\hyperref[II.7.1.2-fr]{7.1.2} et
\hyperref[II.7.1.7-fr]{7.1.7}) un anneau de valuation (resp. un anneau de
valuation discrète) $A$, ayant $L$ pour corps des fractions et dominant
$\mathcal{O}_z$. \emph{A fortiori}, $A$ domine $\mathcal{O}_y$, en posant
$y=p(z)$, et par hypothèse il y a donc un $x\in X$ tel que $A$ domine
$\mathcal{O}_x$. Comme $g$ est propre, la première partie de la démonstration
prouve que $A$ domine aussi un $\mathcal{O}_{x'}$, pour un $x'\in X'$~; il
s'ensuit que $\mathcal{O}_z$ et $\mathcal{O}_{j(x')}=\mathcal{O}_{x'}$ sont
apparentés (\hyperref[I.8.1.4-fr]{I, 8.1.4}), et, comme $P$ est un schéma, cela
entraîne $z=j(x')$ (\hyperref[I.8.2.2-fr]{I, 8.2.2}) et achève la démonstration.
\end{enumerate}

\begin{corollary}[7.3.11]
\label{II.7.3.11-fr}
Soient $X$, $Y$ deux schémas intègres, $f:X\to Y$ un morphisme dominant,
quasi-compact et universellement fermé. Supposons en outre $Y$ affine d'anneau
(intègre) $B$. Alors $\Gamma(X,\mathcal{O}_X)$ est canoniquement isomorphe à
un sous-anneau de la fermeture intégrale de $B$ dans $R(X)$.
\end{corollary}

En effet (\hyperref[I.8.2.1.1-fr]{I, 8.2.1.1}),
$\Gamma(X,\mathcal{O}_X)$ s'identifie à l'intersection des $\mathcal{O}_x$ pour
$x\in X$~; en vertu de (\hyperref[II.7.3.10-fr]{7.3.10}), de
(\hyperref[II.7.1.2-fr]{7.1.2}) et (\hyperref[II.7.1.3-fr]{7.1.3}),
$\Gamma(X,\mathcal{O}_X)$ est par suite contenu dans l'intersection des
anneaux de valuation contenant $B$ et ayant $R(X)$ pour corps des fractions~;
la conclusion résulte alors de (\hyperref[II.7.1.3-fr]{7.1.3}).

\begin{remark}[7.3.12]
\label{II.7.3.12-fr}
Sous les hypothèses de (\hyperref[II.7.3.11-fr]{7.3.11}), et lorsqu'on suppose
que $R(X)$ est une extension de type fini de $R(Y)$, on pourra dans de nombreux
cas conclure que $\Gamma(X,\mathcal{O}_X)$ est un \emph{module de type fini}
sur l'anneau $B=\Gamma(Y,\mathcal{O}_Y)$. Ce sera le cas par exemple lorsque
$B$ est une \emph{algèbre de type fini sur un corps}, car on sait alors que la
fermeture intégrale de $B$ dans une extension de type fini de son corps des
fractions est un $B$-module de type fini ([13], t.~I, p.~267, th.~9)~; la
conclusion résulte alors de (\hyperref[II.7.3.11-fr]{7.3.11}) et du fait que
$B$ est noethérien.

En particulier, \emph{un schéma $X$ propre et affine sur un corps $K$ est
fini}. En effet, en vertu de (\hyperref[II.1.6.4-fr]{1.6.4}),
(\hyperref[II.5.4.6-fr]{5.4.6}) et
(\hyperref[I.6.4.4-fr]{I, 6.4.4, c)}), on peut se borner au cas où $X$ est
\emph{réduit}. En outre, il suffira de prouver que chacun des sous-préschémas
fermés de $X$ ayant pour espace sous-jacent une composante irréductible de $X$
(en nombre fini) est fini sur $K$, si bien que (tenant compte de
(\hyperref[II.5.4.5-fr]{5.4.5})) on est finalement ramené au cas où $X$ est
\emph{intègre}. Mais alors le résultat découle des remarques faites ci-dessus.

Au chapitre III, nous retrouverons cette dernière proposition par d'autres
méthodes et comme conséquence de résultats plus généraux, montrant que si
$f:X\to Y$ est propre et $Y$ localement noethérien, $f_*(\mathcal{F})$ est
cohérent pour tout $\mathcal{O}_X$-Module cohérent $\mathcal{F}$
(\hyperref[III.4.4.2-fr]{III, 4.4.2}).

Notons enfin que le critère (\hyperref[II.7.3.10-fr]{7.3.10}) est pris comme
\emph{définition} des morphismes propres en Géométrie algébrique classique.
Nous ne l'avons mentionné que pour cette raison, le critère
(\hyperref[II.7.3.8-fr]{7.3.8}) semblant plus maniable dans toutes les
applications à notre connaissance.
\end{remark}

